% P17-D5: ordinal aggregation, Arrow's theorem and binary majority.
\chapter{Social Choice and Preference Aggregation}
\label{sec:p17-social-choice}

Individually consistent rankings need not give a consistent collective
ranking when each pair is decided separately. A collective rule must specify
its inputs, its output and the properties it preserves across different
profiles. This chapter develops the aggregation model and complete Arrow
proof in Rubinstein's Lecture 9, followed by a solution of its binary-majority
characterization exercise
\cite[2015 revision, pp.~122--128]{Rubinstein2015LectureNotes}.
The source attributes its four-step Arrow proof to Geanakopolos
\cite[p.~127]{Rubinstein2015LectureNotes}.

The input contains ordinal comparisons. It does not contain interpersonal
comparisons of preference intensity. A theorem about this model therefore
does not decide which ethical objectives a community ought to adopt.

\section{Profiles, rankings and pairwise votes}
\label{sec:p17-social-profiles}

\begin{hbdefinition}{def:p17-social-profile}{A finite electorate and unrestricted ordinal profiles}
Let \(N=\{1,\ldots,n\}\), where \(n\geq1\), be the electorate.
Let \(X\) be a set of alternatives. A \emph{strict ordering} on \(X\)
is an irreflexive, transitive relation that compares every two distinct
alternatives in exactly one direction. Write \(\mathcal L(X)\) for the
set of these orderings. A \emph{profile} is
\(P=(\succ_1,\ldots,\succ_n)\in\mathcal L(X)^N\).
The \emph{unrestricted domain} consists of every such profile.

Let \(\mathcal W(X)\) denote the complete, transitive weak relations on
\(X\), with completeness including equal alternatives. This uses the
same weak/strict conventions as Definition~\ref{def:p17-weak-preference},
now permitting an infinite \(X\). A \emph{social welfare function} is a map
\[
 F:\mathcal L(X)^N\longrightarrow\mathcal W(X).
\]
Write \(\succeq_P=F(P)\), with strict part \(\succ_P\) and indifference
part \(\sim_P\). Social indifference between distinct alternatives is
allowed, even though each input ordering is strict.

A \emph{social choice function}, by contrast, selects one member of \(X\)
at each profile. A social ranking and a selected alternative have different
types. Selecting a maximal alternative from a ranking also requires a
nonempty set of maxima and a rule for choosing among ties.
\index{social welfare function}\index{social choice function}
\index{preference!profile}\index{unrestricted domain}
\end{hbdefinition}

\begin{hbdefinition}{def:p17-pairwise-majority}{Pairwise majority}
For distinct \(x,y\in X\), put
\[
 V_P(x,y)=\#\{i\in N:x\succ_i y\}.
\]
Define \(x\succeq^M_Py\) when \(V_P(x,y)\geq V_P(y,x)\), and
include \(x\succeq^M_Px\). Strict majority means
\(V_P(x,y)>V_P(y,x)\). Equal vote counts give social indifference.
The definition produces a complete relation; transitivity is a question
to be checked, rather than part of the vote-counting operation.
\index{majority rule}\index{voting cycle}
\end{hbdefinition}

\begin{hbtheorem}{thm:p17-social-majority}{proved=annex:p17.social-majority}{The voting cycle and the binary exception}
Pairwise majority is complete and follows every unanimous strict comparison.
Its comparison of a pair depends only on the individual comparisons of
that pair. Nevertheless it need not be transitive.

For three voters and three alternatives, take
\[
 \begin{array}{c|ccc}
 \text{voter}&\text{first}&\text{second}&\text{third}\\ \hline
 1&a&b&c\\
 2&b&c&a\\
 3&c&a&b
 \end{array}
\]
Then \(a\succ^M_Pb\), \(b\succ^M_Pc\) and \(c\succ^M_Pa\), each by
two votes to one. Thus the rule is not a social welfare function on this
unrestricted domain. With exactly two alternatives it is a social welfare
function. With three voters in that binary domain no voter dictates every
social comparison.
\end{hbtheorem}

\section{Aggregation rules and their different commitments}
\label{sec:p17-social-rules}

\begin{hbdefinition}{def:p17-social-aggregation-rules}{Six specified aggregation procedures and rank scores}
The following rules make explicit the examples in the source
\cite[pp.~123--125]{Rubinstein2015LectureNotes}.
Where ranks or finite counts of alternatives are used, assume \(X\) is
finite and nonempty, with \(\ell=\#X\).
\begin{enumerate}
\item A \emph{constant rule} always returns a fixed member of \(\mathcal W(X)\).
\item The \emph{majority-victory score} is
      \(C_P(x)=\#\{z\in X\setminus\{x\}:V_P(x,z)>V_P(z,x)\}\).
      Rank alternatives by decreasing \(C_P\), allowing equal scores.
\item On \(X=\{a,b\}\), a specified biased threshold rule returns
      \(b\succ_Pa\) when \(V_P(b,a)\geq\lceil2n/3\rceil\), and returns
      \(a\succ_Pb\) otherwise. This fixes a strict-output convention for
      the source's two-thirds example; the favored default is \(a\).
\item For a fixed voter \(j\), an \emph{anti-dictatorial rule} reverses
      \(\succ_j\), then adds the diagonal to obtain a weak relation.
\item For a proposed strict ordering \(R\), let
      \[
       D(R,P)=\#\{(x,y,i):x\succ_i y\text{ and }y\,R\,x\}.
      \]
      A \emph{minimum-disagreement rule} chooses an ordering minimizing
      \(D(R,P)\), with ties resolved by a fixed ordering of the finite
      candidate set \(\mathcal L(X)\), and returns its weak closure.
\item A \emph{most-common-ranking rule} chooses an input ordering occurring
      most often, using a fixed ordering of \(\mathcal L(X)\) to resolve
      ties, and returns its weak closure.
\end{enumerate}
Also fix real weights \(w_1>\cdots>w_\ell\). If \(r_i(x)\) is the
position of \(x\) in voter \(i\)'s ordering, the \emph{rank-score rule}
uses \(B_P(x)=\sum_{i=1}^n w_{r_i(x)}\) and ranks by decreasing score.
The three-alternative scores \(2,1,0\) will be used below.
\index{voting!rank scores}\index{voting!minimum disagreement}
\end{hbdefinition}

\begin{hbtheorem}{thm:p17-social-aggregation-rules}{proved=annex:p17.social-aggregation-rules}{These procedures return rankings, but pairwise information can be lost}
Each specified procedure returns a complete, transitive social relation
on its stated domain. Majority-victory and strictly decreasing rank-score
rules follow every unanimous strict comparison.

The majority-victory rule can change its social comparison of \(a,b\)
while every voter's comparison of \(a,b\) remains fixed. In the cyclic
profile above all three victory scores are one. Replacing only voter 2's
ordering by \(b\succ_2a\succ_2c\) gives
\[
 C_{P'}(a)=2,\qquad C_{P'}(b)=1,\qquad C_{P'}(c)=0.
\]
Every individual comparison of \(a,b\) is unchanged, while the social
tie between them becomes \(a\succ_{P'}b\).
\end{hbtheorem}

\begin{hbdefinition}{def:p17-arrow-axioms}{Pareto, independence and dictatorship}
A social welfare function \(F\) satisfies \emph{Pareto} when
\[
 (\forall i\in N,\ x\succ_i y)\quad\Longrightarrow\quad x\succ_Py
\]
for every profile and every distinct pair.
It satisfies \emph{independence of irrelevant alternatives} (IIA) when,
for all profiles \(P,P'\) and distinct \(x,y\),
\[
 \begin{split}
 &(\forall i\in N,\quad x\succ_i y\ \Longleftrightarrow\ x\succ'_i y)\\
 &\hspace{20mm}\Longrightarrow
       (x\succeq_Py\ \Longleftrightarrow\ x\succeq_{P'}y).
 \end{split}
\]
Applying this condition to both ordered pairs preserves both weak
comparisons, hence strict comparison and indifference as well.
IIA compares profiles; it is not a statement that deleting an alternative
from a choice menu leaves every choice unchanged.

The rule is \emph{dictatorial} if there is one fixed \(j\in N\) such that,
at every profile and for every distinct \(x,y\),
\[
                    x\succ_jy\quad\Longrightarrow\quad x\succ_Py.
\]
Because the input is a strict ordering, this determines the entire social
weak relation: it is the reflexive closure of that voter's ordering.
The voter is fixed across profiles and pairs. Agreement with one voter
at a single observed profile does not establish dictatorship.
\index{Pareto condition}\index{independence of irrelevant alternatives}
\index{dictatorship!social choice}
\end{hbdefinition}

\section{Arrow's theorem}
\label{sec:p17-arrow-theorem}

\begin{hblemma}{thm:p17-social-extreme-alternative}{proved=annex:p17.social-extreme-alternative}{An alternative extreme for every voter is socially extreme}
Assume \(X\) contains at least three alternatives and \(F\) satisfies
Pareto and IIA. Fix \(b\in X\). At a profile where every voter places
\(b\) either strictly above all other alternatives or strictly below all
of them, the social relation also places \(b\) strictly above all the
others or strictly below all the others. In particular it cannot tie with
another alternative at such a profile.
\end{hblemma}

\begin{hbtheorem}{thm:p17-social-arrow}{proved=annex:p17.social-arrow}{Arrow's theorem for a finite electorate}
Let \(N=\{1,\ldots,n\}\) with \(n\geq1\), and let \(X\) contain at
least three alternatives. Every social welfare function on the unrestricted
strict-profile domain \(\mathcal L(X)^N\) satisfying Pareto and IIA is
dictatorial.

The electorate is finite. The alternative set need not be finite.
Social outputs are complete and transitive and may initially be specified
to allow ties; dictatorship rules out ties between distinct alternatives
on this strict-input domain. Equivalently, no nondictatorial function
satisfies all these requirements
\cite[pp.~125--126]{Rubinstein2015LectureNotes}.
\end{hbtheorem}

\begin{hbmethod}{def:p17-arrow-pivot}{Finding the voter used in the proof}
Choose a reference alternative \(b\) and any profile on the remaining
alternatives. For \(m=0,\ldots,n\), form \(P^m\) by placing \(b\) at
the top for voters \(i\leq m\), and at the bottom for voters \(i>m\).
The extreme-alternative lemma makes \(b\) socially top or bottom at
every step. Pareto makes it bottom at step zero and top at step \(n\).
Let \(m^*\) be the first step at which it is socially top.

The proof uses IIA to show that the same \(m^*\) works for any rankings
of pairs not involving \(b\). It then inserts \(b\) between two alternatives
for voter \(m^*\), and at opposite extremes for the other voters.
Comparisons with \(P^{m^*-1}\) and \(P^{m^*}\) force agreement with
voter \(m^*\) on that pair. A third alternative transfers the conclusion
to pairs involving \(b\), in both directions. The annex supplies all
these steps, including the possible weak-comparison cases.

This is a proof using the function's behavior on a whole unrestricted
domain. It is not a statistical method for inferring dictatorship from
a short list of observed votes.
\end{hbmethod}

\begin{hbtheorem}{thm:p17-social-arrow-boundaries}{proved=annex:p17.social-arrow-boundaries}{Why the hypotheses matter}
Each of the following is a counterexample to dropping a stated requirement:
\begin{enumerate}
\item Without Pareto, reversing one voter's ordering gives a nondictatorial
      social welfare function satisfying IIA.
\item Without IIA, finite rank scores give a Pareto social welfare function
      that need not be dictatorial. With two voters, compare
      \[
       P=(a>b>c,\ b>c>a),\qquad
       P'=(a>c>b,\ b>a>c).
      \]
      Using scores \(2,1,0\), the scores of \((a,b,c)\) are
      \((2,3,1)\) at \(P\) and \((3,2,1)\) at \(P'\).
      The individual \(a,b\) comparisons remain unchanged.
\item Without social transitivity, the three-voter majority rule has
      Pareto and IIA but has the displayed voting cycle.
\item With only two alternatives, three-voter majority is a nondictatorial
      Pareto and IIA social welfare function.
\end{enumerate}
The notation \(a>b>c\) in the two profiles means the voter's strict
ordering, not an independent numerical order on the alternatives.
\end{hbtheorem}

The unrestricted profile domain, completeness of the social output and
finite electorate are also premises of the theorem. The four examples
above do not characterize every way of changing those premises.
For example, restricting admissible profiles requires a separate theorem
about that restricted domain.

\section{Two alternatives with individual indifference}
\label{sec:p17-binary-majority-characterization}

\begin{hbdefinition}{def:p17-binary-votes}{Signed binary preferences and three symmetry conditions}
Now let \(X=\{a,b\}\) and allow individual indifference. Encode a voter's
weak preference by \(s_i=1\) for \(a\succ_i b\), \(s_i=0\) for
\(a\sim_i b\), and \(s_i=-1\) for \(b\succ_i a\).
A binary aggregation rule is a function
\[
                   f:\{-1,0,1\}^n\longrightarrow\{-1,0,1\},
\]
with the same interpretation for its output. This is a new input domain,
distinct from the strict-input domain of Arrow's theorem.

\emph{Anonymity} means that permuting voters does not change \(f\).
\emph{Neutrality} means \(f(-s)=-f(s)\).
\emph{Positive responsiveness} means that whenever one coordinate
increases by one, with every other coordinate fixed, then
\[
                   f(s)\geq0\quad\Longrightarrow\quad f(s')=1.
\]
The permitted increases are \(-1\) to \(0\) and \(0\) to \(1\).
Symmetrically, a decrease by one with \(f(s)\leq0\) must give
\(f(s')=-1\). Under neutrality this second direction follows from the
first; it is included in the condition itself when discussing independence.
Thus improving support for an alternative turns a social tie in its favor,
and preserves an already strict social preference for it
\cite[p.~128, problem~1]{Rubinstein2015LectureNotes}.
\index{anonymity!voting}\index{neutrality!voting}
\index{positive responsiveness}
\end{hbdefinition}

\begin{hbtheorem}{thm:p17-social-may}{proved=annex:p17.social-may}{Binary majority characterized by anonymity, neutrality and responsiveness}
For each \(n\geq1\), exactly one binary rule has all three properties:
\[
 f(s)=\operatorname{sgn}\left(\sum_{i=1}^n s_i\right),
\]
where \(\operatorname{sgn}(t)\) is \(1,0,-1\) according as the real
number \(t\) is positive, zero or negative.
In particular equal numbers of strict supporters on the two sides produce
social indifference; indifferent voters contribute zero.
This is the binary characterization requested in the source's exercise.
\end{hbtheorem}

\begin{hbexample}{def:p17-binary-count-example}{A tied electorate and a decisive change}
For \(s=(1,1,-1,-1,0)\), the vote sum is zero, so the characterized
rule returns indifference. Changing the last voter from \(0\) to \(1\)
makes the sum one and the social comparison strictly favor \(a\).
Changing one opposing voter from \(-1\) to \(0\) has the same effect.
The rule responds to the count difference; the fraction of all voters
strictly supporting \(a\) need not exceed one half when some voters
are indifferent.
\end{hbexample}

\begin{hbtheorem}{thm:p17-social-may-boundaries}{proved=annex:p17.social-may-boundaries}{Independence of the three binary conditions}
For two voters the conditions are independent:
\begin{enumerate}
\item \(f(s)=\operatorname{sgn}(2s_1+s_2)\) is neutral and positively
      responsive, but is not anonymous.
\item The constant rule \(f(s)=1\) is anonymous and positively responsive,
      but is not neutral.
\item The constant rule \(f(s)=0\) is anonymous and neutral, but is not
      positively responsive.
\end{enumerate}
\end{hbtheorem}

\begin{proofboundary}{def:p17-social-choice-next}{Further collective-choice regimes}
All results asserted in this chapter have full annex proofs.
Chapter~\ref{sec:p17-single-peaked-majority} supplies the single-peaked
domain, full majority-order and median-winner results, including the
source's Problem~2. Further required treatments include other preference
domains, aggregation and voting methods, strategic misreporting and
the precise domain/range assumptions of social-choice impossibility
theorems. Arrow's theorem here concerns social rankings and IIA; it does
not by itself prove a theorem about strategic manipulation of a function
selecting one alternative. The source's other related-results discussion
and Problems 3--5 require their own statements and constructions.

The matching comparison in \cref{thm:p18-marriage-argumentation,thm:p18-deferred-acceptance,thm:p18-matching-welfare-boundary}
distinguishes stability, the proposing side's optimality and a specified
aggregate-welfare objective. \Cref{thm:p18-five-person-pairing} also
separates maximal admissibility from maximizing the number matched.
The coalitional core and NM solution comparison in
\cref{thm:p18-imputation-core,thm:p18-nm-argumentation}
uses coalition values and imputations; it is not a Nash-equilibrium
claim about strategic action profiles.
Finite strategic profiles, both original Nash existence constructions
and the minimax and welfare comparisons are supplied in
\cref{sec:p17-finite-strategic-games}. The precise solution-object
comparison is \cref{thm:p17-strategic-welfare-boundaries}.
\end{proofboundary}
